\section{Related Work}
% GOAL: 600-800 words. Organized by 4 thematic groups. Compare and contrast. Position EvoBench clearly.

% === GROUP 1: Repository-level coding benchmarks ===
% SWE-bench family: execution-grounded, real GitHub issues, patch verification
% Key insight to convey: EvoBench inherits execution-grounding but extends to serial dependency

\textbf{Repository-level coding benchmarks.}
SWE-bench~\cite{jimenez2024swebench} established the paradigm of evaluating agents on real GitHub issues with execution-based patch verification, demonstrating that repository-level context and test suites provide more realistic assessment than isolated function synthesis~\cite{chen2021codexhumaneval}.
SWE-bench Verified~\cite{chowdhury2024swebenchverified} refined this with human-validated task specifications, and subsequent extensions have broadened coverage to multilingual settings~\cite{zan2025swebenchmultilingual}, longer-horizon tasks~\cite{deng2025swebenchpro}, mobile development~\cite{yang2025swebenchmobile}, and hardware engineering~\cite{hwbench2026}.
FEA-Bench~\cite{feabench2025} extended the paradigm to feature implementation rather than bug fixing.
These benchmarks share a core property: each task is \emph{independent}---the agent starts from a clean repository state and produces a single patch.
\eva{} inherits the execution-grounded verification tradition but fundamentally differs by requiring agents to progress along a \emph{serial dependency chain} where intermediate artifacts accumulate across tasks.

% === GROUP 2: Interactive / tool-use agent benchmarks ===
% AgentBench, GAIA, OSWorld, Terminal-Bench, WebArena
% Key insight: these cover breadth but not serial depth

\textbf{Interactive agent benchmarks.}
AgentBench~\cite{liu2024agentbench} evaluates LLM agents across eight interactive environments spanning OS operations, databases, and web tasks.
GAIA~\cite{mialon2024gaia} tests general assistant competence on multi-step reasoning questions requiring tool use.
OSWorld~\cite{xie2024osworld} and WebArena~\cite{zhou2024webarena} assess GUI and web navigation in real computer environments.
Terminal-Bench~\cite{merrill2026terminalbench} provides execution-grounded evaluation in command-line settings.
These benchmarks emphasize \emph{breadth} across diverse interaction surfaces.
\eva{} instead emphasizes \emph{depth} along a single dependency chain, enabling precise measurement of cascading failure, intermediate artifact continuity, and process cost accumulation that broader suites cannot isolate.

% === GROUP 3: Paired / augmentation-oriented evaluation ===
% SkillsBench, SWE-Skills-Bench
% Key insight: paired evaluation methodology is directly relevant to our resurrection design

\textbf{Paired and augmentation-oriented evaluation.}
SkillsBench~\cite{li2026skillsbench} introduced the methodology of evaluating every task under both vanilla and skill-augmented conditions, demonstrating that augmentation efficacy is context-dependent and requires paired measurement.
SWE-Skills-Bench~\cite{sweskillsbench2026} applied this methodology to software engineering, finding that most public skills provide negligible benefit.
\eva{} adopts the paired evaluation philosophy but applies it to a different axis: \emph{with vs.\ without resurrection}, measuring how much diagnostic signal is recovered when predecessor failures are bypassed.
This parallels SkillsBench's ``no-skill vs.\ skill'' comparison but addresses the orthogonal question of \emph{failure decomposition} rather than \emph{augmentation efficacy}.

% === GROUP 4: Self-evolving / long-horizon agents ===
% Frontier-Eng, RE-Bench, Voyager, continual learning
% Key insight: EvoBench provides a grounded testbed for studying evolution behavior

\textbf{Long-horizon and self-evolving agent evaluation.}
Recent work has begun examining agents over extended trajectories.
Frontier-Eng~\cite{frontiereng2026} benchmarks self-evolving agents on real-world R&D workflows.
RE-Bench~\cite{rebench2024} and PaperBench~\cite{paperbench2025} evaluate frontier research engineering tasks.
Voyager~\cite{wang2023voyager} demonstrated skill accumulation in Minecraft environments.
These works highlight the importance of long-horizon evaluation but lack \emph{controlled intermediate state injection}---when an agent fails early, the entire trajectory is typically discarded.
\eva{}'s resurrection mechanism provides a principled way to continue evaluation after early failure, enabling measurement of both evolutionary gain (using self-produced artifacts) and latent downstream capability (using golden artifacts).

% === POSITIONING SUMMARY ===
Table~\ref{tab:benchmark_comparison} positions \eva{} against existing benchmarks across five evaluation dimensions.

\begin{table}[t]
\centering
\small
\begin{tabular}{lcccccc}
\toprule
\textbf{Benchmark} & \textbf{Serial} & \textbf{Exec.} & \textbf{Resurr.} & \textbf{Process} & \textbf{Artifact} & \textbf{Tasks} \\
\midrule
SWE-bench~\cite{jimenez2024swebench} & --- & \checkmark & --- & --- & --- & 2.3K \\
AgentBench~\cite{liu2024agentbench} & --- & \checkmark & --- & --- & --- & 8 envs \\
GAIA~\cite{mialon2024gaia} & --- & --- & --- & --- & --- & 466 \\
OSWorld~\cite{xie2024osworld} & --- & \checkmark & --- & --- & --- & 369 \\
WebArena~\cite{zhou2024webarena} & --- & \checkmark & --- & --- & --- & 812 \\
Terminal-Bench~\cite{merrill2026terminalbench} & --- & \checkmark & --- & --- & --- & 50 \\
SkillsBench~\cite{li2026skillsbench} & --- & \checkmark & paired & --- & --- & 1.5K \\
RE-Bench~\cite{rebench2024} & --- & \checkmark & --- & partial & --- & 7 \\
\midrule
\eva{} & \checkmark & \checkmark & \checkmark & \checkmark & \checkmark & 6 \\
\bottomrule
\end{tabular}
\caption{Benchmark comparison. \checkmark = supported, --- = not supported. \emph{Serial} = tasks form a dependency chain; \emph{Exec.} = execution-grounded scoring; \emph{Resurr.} = resurrection mechanism; \emph{Process} = process metrics as first-class output; \emph{Artifact} = intermediate artifact continuity tracked.}
\label{tab:benchmark_comparison}
\end{table}

\eva{} is the only benchmark that combines serial dependency, execution-grounded scoring, resurrection, process metrics, and artifact continuity tracking.
While SWE-bench and SkillsBench share execution-grounding and paired evaluation (respectively), neither addresses serial dependency or intermediate artifact flow.
RE-Bench records partial process metrics but lacks resurrection and artifact tracking.
The combination of these five dimensions enables \eva{} to measure evaluation axes that no existing benchmark captures simultaneously.